\chapter{Kryptologie}\label{ch:Kryptologie}

\lstset{style=TigerJython, numbers=none}


\section{Zeichenketten in Python}
Im Nachfolgenden schauen wir uns einige Python-Befehle an, die dazu dienen, mit Zeichenketten ($=$ engl. \textit{strings}) zu arbeiten.
Strings werden wir insbesondere benötigen, um eigene Python-Programme zu schreiben, mit denen wir Texte verschlüsseln und entschlüsseln können. 
Genauer gesagt bezeichnen wir Texte in Python als Zeichenketten. 
Zeichenketten werden in Python innerhalb von einfachen oder doppelten Anführungszeichen geschrieben, also entweder \lstinline|'...'| oder \lstinline|"..."|. 
Innerhalb der Anführungszeichen können beliebige Zeichen stehen, beispielsweise Buchstaben, Leerzeichen, Spezialzeichen oder Zahlen.

\begin{mychallenge}
Was könnte der Nutzen davon sein, dass es zwei Möglichkeiten gibt und nicht nur eine, Zeichenketten zu schreiben (\lstinline|'...'| oder \lstinline|"..."|)? 
\begin{myanswer}
Eine Möglichkeit, die uns durch die unterschiedlichen Schreibweisen geboten wird, ist, Anführungszeichen auszugeben. Beispielsweise könnten Sie so einen Text, der Anführungszeichen beinhaltet, auf der Konsole drucken:
\begin{lstlisting}
print('Meine Freunde nennen mich "Mr. X"')
\end{lstlisting}
\end{myanswer}
\end{mychallenge}

Mit der \textbf{Länge} von Zeichenketten ist die Anzahl der Zeichen zwischen den Anführungszeichen gemeint: Die Zeichenkette \lstinline|"Hello World"| hat beispielsweise eine Länge von 11 Zeichen. Die Länge einer Zeichenkette kann mit dem Ausdruck \pythoninline{len(Klartext)} bestimmt werden.

\begin{myexercise}
  Bestimmen Sie die Länge Ihres vollen Namens mithilfe des Befehls \pythoninline{len()} in Python.
\end{myexercise}

Zeichenketten können verkettet, also aneinandergehängt werden mit dem Befehl \lstinline|"Text1" + "Text2" + "Text3"| usw. 
Falls Sie eine Zeichenkette mehrfach drucken wollen, können Sie diesen mit einer Zahl multiplizieren, die dann die Anzahl Wiederholungen der Zeichenkette bestimmt. So ist der Ausdruck \lstinline|"a" * 3| beispielsweise gleichbedeutend mit einer Zeichenkette \lstinline|"aaa"|.

\begin{myexercise}\label{ex:pfeil}
Führen Sie das nachfolgende Programm aus und erklären Sie, was das Programm tut.
  \begin{lstlisting}
    print(" " * 3 + "X" * 1)
    print(" " * 2 + "X" * 3)
    print(" " * 1 + "X" * 5)
    print(" " * 0 + "X" * 7)
    print(" " * 3 + "X" * 1)
    print(" " * 3 + "X" * 1)
  \end{lstlisting}
\end{myexercise}

\begin{myexercise}
Erstellen Sie nun ein Programm, das ein Herz zeichnet, ähnlich zu \cref{ex:pfeil}, also lediglich mit Zeichenketten-Ausdrücken und dem Befehl \pythoninline{print()}.
\begin{myanswer}
\begin{lstlisting}
print(" " * 2 + "X" * 1 + " "*4 + "X" * 1 + " "*1)
print(" " * 1 + "X" * 8 + " "*1)
print(" " * 2 + "X" * 6 + " "*1)
print(" " * 3 + "X" * 4 + " "*2)
print(" " * 4 + "X" * 2 + " "*2)
\end{lstlisting}
\end{myanswer}
\end{myexercise}


\subsection{Zeichen aus Zeichenketten auslesen}
Im Nachfolgenden schauen wir uns einige Übungen an, mit denen wir die einzelnen Zeichen aus Zeichenketten herauslesen können.

\begin{myexercise}
  Führen Sie das nachfolgende Programm aus und erklären Sie, was das Programm tut.
  \begin{lstlisting}
    text = "EASY"
    for buchstabe in text:
      print(buchstabe)
  \end{lstlisting}
  \begin{myanswer}
    \begin{lstlisting}
      E
      A
      S
      Y
    \end{lstlisting}
    \pythoninline{for x in mystring} ist eine Schleife über die Zeichenkette \pythoninline{mystring}. Dabei nimmt \pythoninline{x} nacheinander jeden Buchstaben in dieser Zeichenkette genau einmal an.
  \end{myanswer}
\end{myexercise}

\begin{myexercise}\label{myexercise:index}
  Führen Sie das nachfolgende Programm aus und erklären Sie, was das Programm tut.
  \begin{lstlisting}
    Klartext = "Schweiz"
    print(Klartext[0])
    print(Klartext[1])
    print(Klartext[2])
    print(Klartext[3])
    print(Klartext[4])
    print(Klartext[5])
    print(Klartext[6])
  \end{lstlisting}
  \begin{myanswer}
    \begin{lstlisting}
      S
      c
      h
      w
      e
      i
      z
    \end{lstlisting}
    \pythoninline{mystring[0]} ist der erste Buchstabe der Zeichenkette \pythoninline{mystring}, \pythoninline{mystring[1]} der zweite und so weiter.
  \end{myanswer}
\end{myexercise}

\begin{myexercise}
  Führen Sie das nachfolgende Programm aus und erklären Sie, was das Programm tut.
  \begin{lstlisting}
    Klartext = "Schweiz"
    for i in range(len(Klartext)):
      print(Klartext[i])
  \end{lstlisting}
  \begin{myanswer}
    Dieses Programm ist äquivalent zum Programm in \cref{myexercise:index}. 
    Hier wird allerdings eine \pythoninline{for}-Schleife über die Länge des Texts verwendet.
  \end{myanswer}
\end{myexercise}

Der Befehl \pythoninline{for i in range(len(Klartext))} erstellt eine Variable \pythoninline{i} innerhalb der \pythoninline{for}-Schleife, welche jedes Zeichen von 0 bis zur Länge von \pythoninline{Klartext} minus 1 geht. Weshalb minus 1? Das erste Zeichen von \pythoninline{Klartext} wird in Python mit \pythoninline{Klartext[0]} ausgelesen, das letzte mit \pythoninline{Klartext[len(Klartext)-1]}, da wir bei 0 zu zählen beginnen und nicht bei 1. Mit dem Ausdruck \pythoninline{Klartext[i]} wird also das \pythoninline{i}-te Zeichen der Zeichenkette \pythoninline{Klartext} ausgelesen, wobei \pythoninline{i} von 0 bis zu (Länge des Klartexts minus 1) geht. 

Der Befehl \pythoninline{for i in range(len(Klartext))} kann auch geschrieben werden also Befehl \pythoninline{for i in range(0, len(Klartext), 1)}, wobei dies meint:
\begin{itemize}
\item \pythoninline{i} beginnt bei 0
\item \pythoninline{i} geht bis zu \pythoninline{len(Klartext)-1}
\item \pythoninline{i} vergrössert sich in 1er-Schritten
\end{itemize}
Dieser Befehl könnte auch verwendet werden, um bei einer beliebigen Zahl zu starten (nicht notwendigerweise 0), und um \pythoninline{i} in Schritten grösser als 1 zu vergrössern. Die allgemeine Syntax des Befehls lässt sich also wie folgt zusammenfassen: \pythoninline{for i in range(start, ende, schrittgroesse)}.

\begin{mychallenge}
  Eine Person verrät uns lediglich die Vorwahl ihrer $10$-stelligen Mobiltelefonnummer. Zudem verrät sie Ihnen auch, dass ihre Telefonnummer gerade ist (also auf 2, 4, 6, 8 oder 0 endet).
  Damit gibt es nur noch $5$ Millionen Kombinationen, die es (im schlimmsten Fall) auszuprobieren gilt. 
  Schreiben Sie ein Python-Programm, welches die $2000$ kleinsten, geraden Nummern auflistet. 
  Die erste Nummer sollte $0790000000$ sein, die letzte $0790001998$. 
  Das Programm soll die Nummern als Zeichenkette (eine Nummer pro Zeile) ausgeben.
  
  Tipps:
  \begin{itemize}
    \item Mit \pythoninline{str(num)} 
      wird aus der Zahl \pythoninline{num} eine Zeichenkette.
      Beispielsweise gibt uns \pythoninline{str(15)}
      die Zeichenkette \lstinline[language=Python,basicstyle=\ttfamily\normalsize]!"15"!.
    \item Erinnern Sie sich, was die Operation \pythoninline{+} in dem Ausdruck \lstinline[language=Python,basicstyle=\ttfamily\normalsize]!"In" + "form" + "atik"! macht?
  \end{itemize}
  \begin{myanswer}
    \begin{lstlisting}
      for num in range(790000000, 790002000,2):
        print("0" + str(num)) # füge die Zeichenkette "0" vorne an
    \end{lstlisting}
  \end{myanswer}
\end{mychallenge}

\section{Alphabete, Wörter, Sprachen}
\begin{mydefinition}{\textbf{Alphabet}}\label{mydefinition:Alphabet}

Eine endliche nichtleere Menge $\Sigma$ heisst \textit{Alphabet}. Die \textit{Elemente eines Alphabets} werden \textit{Buchstaben} (\textit{Zeichen, Symbole}) genannt.
\end{mydefinition}
\cref{mydefinition:Alphabet} entspricht unserer intuitiven Vorstellung eines Alphabets: 
\begin{itemize}
  \item Um Text darstellen zu können, muss ein Alphabet mindestens ein Symbol enthalten ($\to$ nichtleer).
  \item Damit man sich auf einen Satz von Zeichen einigen kann, darf das Alphabet nicht unendlich gross sein ($\to$ endlich).
\end{itemize}

\clearpage
\noindent
Wir listen nun einige der Alphabete auf, die in der Mathematik und Informatik häufig verwendet werden:
\begin{itemize}
    \item $\Sigma_{\text{bool}} = \lrc{0,1}$ ist das Boole'sche Alphabet, mit dem digitale Rechner arbeiten.
    \item $\Sigma_{\text{lat}} = \lrc{a,b,c,...,z,A,B,...,Z}$ ist das lateinische Alphabet.
    \item $\Sigma_{\text{Tastatur}} = \lrc{a,b,c,...,z,A,B,C,...,Z,\text{\textvisiblespace}\,,>,<,(,),...,\#,?,!}$ ist das Alphabet aller Symbole, die mit der englischen Tastatur getippt werden können. Dabei ist \textvisiblespace\ das Symbol für das Leerzeichen / Leersymbol.
    \item $\Sigma_{\text{greek}} = \lrc{\alpha, \beta, \gamma, ..., \omega, A, B, \Gamma, ..., \Omega}$ ist das griechische Alphabet.
    \item $\Sigma_m = \lrc{0,1,2,...,m-1}$ für jedes $m\in\Nunit$ ($\N$ ohne $0$) ist ein Alphabet für die $m$-adische Darstellung von Zahlen.
\end{itemize}
\textit{Wörter} werden wir als endliche Folgen von Buchstaben ansehen.
\begin{mydefinition}{\textbf{Wort}}\label{mydefinition:Wort}
\begin{itemize}
  \item Sei $\Sigma$ ein Alphabet. Ein \textit{Wort} über $\Sigma$ ist eine endliche (möglicherweise leere) Folge von Buchstaben aus $\Sigma$. 
  \item Das \textit{leere Wort} $\lambda$ ist die leere Buchstabenfolge.
  \item Die \textit{Länge} $|w|$ eines Wortes $w$ ist die Länge des Wortes als Folge, das heisst die Anzahl der Vorkommen von Buchstaben in $w$.
\end{itemize}
\end{mydefinition}

\begin{myexample}
  \begin{aenum}
  \item $w = 1,0,0,1,0$ ist ein Wort über dem Alphabet $\Sigma_{\text{bool}}$ und $|w| = 5$, da $w$ eine Folge von $5$ Buchstaben ist.
  \item $u = M,\text{\textvisiblespace}\,, j$ ist ein Wort über $\Sigma_{\text{Tastatur}}$ und $|u| = 3$, da $u$ eine Folge von $3$ Buchstaben ist.
  \item Das leere Wort $\lambda$ ist ein Wort über jedem Alphabet und es gilt $\abs{\lambda} = 0$.
  \end{aenum}
\end{myexample}

\begin{myexercise}
Was ist $\Sigma_{10}$?
\begin{myanswer}
\noindent
\begin{align*}
    \Sigma_{10} = \lrc{0,1,2,...,9}
\end{align*}
ist das Alphabet für die Darstellung von Zahlen im Dezimalsystem.
\end{myanswer}
\end{myexercise}

\begin{mydefinition}{\textbf{Sigma Stern}}

Sei $\Sigma$ ein Alphabet. $\Sigma^*$ (gesprochen: \textit{Sigma Stern}) ist die Menge aller Wörter über $\Sigma$, also die Menge aller endlichen Folgen von Symbolen aus $\Sigma$.
\end{mydefinition}

\begin{myremark}
  \begin{aenum}
  \item Wir werden im Folgenden Wörter ohne Kommas schreiben. Anstelle von $1,0,0,1,0$ werden wir lediglich $10010$ schreiben. Allgemein, werden wir anstelle von $x_1,x_2,...,x_n$ einfach $x_1x_2...x_n$ schreiben.
  \item In der deutschen Sprache existieren die zwei Ausdrücke \enquote{Wörter} und \enquote{Worte}. Dies sind keine Synonyme. \textit{Wörter} bezeichnet den Plural von \textit{Wort} (\enquote{Im Duden stehen viele Wörter}). Der Ausdruck \textit{Worte} bezieht sich auf Gedankenkonstrukte (\enquote{Sie sprach weise Worte}).
  \end{aenum}
\end{myremark}


\section{Motivation}

Alice möchte Bob eine Nachricht zukommen lassen. Diese Nachricht wird \tib{Klartext}\index{Klartext} genannt. 
In der Praxis wird Alice 
\begin{figure}[H]
    \centering
    \begin{tikzpicture}[mylabel/.style={text width=20mm, align=center},scale=0.8, every node/.style={scale=0.8}]
    \node[alice, minimum width=1.5cm, 
        label={[mylabel, anchor=west]north east:{}}, 
        label={[mylabel, anchor=east]north west:{\textcolor{Blue}{Klartext} $\big\downarrow\text{\tiny (Chiffrierung)}$\\
        \textcolor{Red}{Geheimtext}}}] (alice) {Alice (Sender)};
    \node[bob, minimum width=1.5cm, right=5cm of alice, mirrored,
        label={[mylabel, anchor=east]north west:{}}, 
        label={[mylabel, anchor=west]north east:{\textcolor{Red}{Geheimtext} $\big\downarrow\text{\tiny (Dechiffrierung)}$\\
        \textcolor{Blue}{Klartext}}}] (bob) {Bob (Empfänger)};
    
    \draw[->, thick, black] ([yshift=0mm]alice.east) coordinate(aux-a)-- coordinate[near start] (aux) node[above, text width=3cm, align=center]{Übermittlung über unsicheren Kanal e.g. Internet, Telefonie \\
    \textcolor{Red}{Geheimtext}} (aux-a-|bob.west);

    \node[devil, minimum width=1.5cm, below right=13mm and 1.6cm of alice] (eve) {Eve (Gegner / Abhörer)};
    \draw[<-, gray] (eve)--(aux) node[midway,right]{abhören};
    \end{tikzpicture}
    \caption{Dieses Schema zeigt die Kommunikation zwischen Alice (Sender) und Bob (Empfänger).}
    \label{fig:schema_geheimtext}
\end{figure}
\clearpage

\section{Geheimschriften}
\section{Kryptosysteme}
\begin{figure}[H]
    \centering
    \begin{tikzpicture}[mylabel/.style={text width=20mm, align=center},scale=0.8, every node/.style={scale=0.8}]
    \node[alice, minimum width=1.5cm, 
        label={[mylabel, anchor=west]north east:{}}, 
        label={[mylabel, anchor=east]north west:{\textcolor{Blue}{(Klartext, Schlüssel)} $\big\downarrow\text{\tiny (Verschlüsselung)}$\\
        \textcolor{Red}{Kryptotext}}}] (alice) {Alice (Sender)};
    \node[bob, minimum width=1.5cm, right=5cm of alice, mirrored,
        label={[mylabel, anchor=east]north west:{}}, 
        label={[mylabel, anchor=west]north east:{\textcolor{Red}{(Kryptotext, Schlüssel)} $\big\downarrow\text{\tiny (Entschlüsselung)}$\\
        \textcolor{Blue}{Klartext}}}] (bob) {Bob (Empfänger)};
    
    \draw[->, thick, black] ([yshift=0mm]alice.east) coordinate(aux-a)-- coordinate[near start] (aux) node[above, text width=3cm, align=center]{Übermittlung über unsicheren Kanal e.g. Internet, Telefonie \\
    \textcolor{Red}{Kryptotext}} (aux-a-|bob.west);

    \node[devil, minimum width=1.5cm, below right=13mm and 1.6cm of alice] (eve) {Eve (Gegner / Abhörer)};
    \draw[<-, gray] (eve)--(aux) node[midway,right]{abhören};
    \end{tikzpicture}
    \caption{Dieses Schema zeigt die Kommunikation zwischen Alice (Sender) und Bob (Empfänger) unter Verwendung eines Kryptosystems.}
    \label{fig:schema_kryptotext}
\end{figure}

\section{Kryptoanalyse}

\section{Moderne Kryptologie}


